Consideriamo ora le clausole \(\set{A,B},\set{\neg A, \neg B}\). Si potrebbe pensare di applicare la risoluzione su entrambi i clash insieme, ottenendo qunidi \(\square\). Questa cosa però non è possibile, poiché fare questo ci farebbe perdere la correttezza, perchè queste due clausole sono banalmente soddisfacibili. La regola infatti decide prima un letterale su cui fare clash, e poi applica la risoluzione, e non è possibile fare entrambe le cose insieme su più letterali.

\begin{theorem}{}
  \(\textsc{Resolution}(S)\) termina sempre e l'insieme di saturazione \(S_0\) avrà al più \(3^{\abs{\AP[S]}}\) clausole.
\end{theorem}

\begin{proof}
  Sappiamo che \(S\) è un insieme finito di clausole finite. Sia dunque \(\AP[S]\) l'insieme di proposizioni atomiche che compaiono in \(S\). Chiaramente \(\abs{C}\leq\abs{\AP[S]}\). Questo perché sono state tolte le clausole banali e la clausole sono insiemi, quindi ogni clausole può contenere al massimo una versione del letterale per ogni proposizione atomica.

  Possiamo dunque vedere ogni clausola come una funzione che associa a \(P\in \AP[S]\) una delle tre possibilità: \(P\) è presente nella clausola, \(\neg P\) è presente nella clausola, o nessuna delle due è presente nella clausola. Dunque il numero di clausole che è possibile formare con le proposizioni atomiche di \(S\) è al più \(3^{\abs{\AP[S]}}\)
\end{proof}

Abbiamo già visto che \(S\models_R \square \implies S \in \Unsat\), vorremmo dimostrare anche l'altro verso, ovvero che \(S \in \Unsat \implies S \models_R \square\).

La completezza del calcolo di risoluzione garantisce anche la completezza dell'algoritmo di risoluzione, poiché l'algoritmo genera tutto il generabile.

\begin{proof}
  Per induzione su \(\abs{\AP[S]} = n\).

  Se \(n = 0\), allora l'unica clausola che è possibile formare è la clausola vuota \(\square\). Quindi \(S = \set{\square}\), e ovviamente  \(S\models_R\square\).
  Per \(n>0\), abbiamo l'ipotesi induttiva che per ogni \(S'\) con \(\abs{\AP[S']}<n\), \(S'\in\Unsat\implies S'\models_R\square\).
  
  Quello che vogliamo fare è quindi rimuovere in qualche modo una proposizione atomica. Fissiamo dunque \(P\in \AP[S]\).
  Costruiamo allora gli insiemi \(S_P\) e \(S_{\neg P}\) tali che:
  \begin{itemize}
    \item \(P \notin \AP[S_P]\) e \(P\notin \AP[S_{\neg P}]\)
    \item \(\AP[S_P] \subset \AP[S]\) e \(\AP[S_{\neg P}]\subset \AP[S]\)
    \item \(S_P, S_{\neg P} \in \Unsat\)
  \end{itemize}

  Andiamo a definire quindi questi due insiemi come:
  \[S_P = \set{C\setminus P}[C\in S \text{ e } \neg P \in C]\]
  \[S_{\neg P} =  \set{C\setminus \neg P}[C\in S \text{ e }  P \in C]\]

  Ovviamente nessuna clausola nè di \(S_P\), nè di \(S_{\neg P}\) contengono nè \(\neg P\) nè \(P\).
  Dunque, nè \(\AP[S_P]\) nè \(\AP[S_{\neg P}]\) contengono \(P\), quindi \(\AP[S_P] \subset \AP[S]\) e \(\AP[S_{\neg P}]\subset \AP[S]\).

  Per poter applicare l'ipotesi induttiva ci manca solamente garantire che questi due insiemi siano insoddisfacibili.

  Prima di vedere ciò è importante notare che tali insiemi non sono disgiunti. In particolare quelle che sono in entrambi erano già presente in \(S\). Questo perché per stare in entrambi gli insiemi non deve contenere ne \(P\) ne \(\neg P\), e quindi la sottrazione non la cambia.
  Inoltre non sono vuoti. Supponiamo infatti per assurdo che \(S_P=\emptyset\). Dunque ogni clausola in \(S\) contiene \(\neg P\).
  Ma allora basta prendere un assegnamento che assegni \(0\) a \(P\), e tale assegnamento soddisferebbe \(S\), il che contraddice l'ipotesi che \(S\) sia insoddisfacibile.

  Raggioniamo ora sull'insoddisfacibilità di \(S_P\), ovviamente l'altro sarà simile.

  Supponiamo \(S_P\) soddisfacibile per assurdo. Allora esiste un assegnamento \(\alpha:\AP[S_P]\to\set{0,1}\), tale che \(\Val{\alpha}[S_p] = 1\). Definiamo allora \(\alpha^p\) come:

  \[\alpha^P(X)\begin{cases}
    0 &, X = P\\
    \alpha(X) &, X \neq P
  \end{cases}\]

  Ovviamente questo è un assegnamento per \(S\), poiché aggingendo \(P\) a \(\AP[S_P]\) si ottiene \(\AP[S]\).

  In particolare possiamo mostrare che \(\Val{\alpha^P}[S] = 1\). Questo è vero se e solo se \(\forall C \in S: \Val{\alpha^P}[C] =1\).
  Ora, preso un \(C\) qualsaisi, se \(\neg P \in C \implies \Val{\alpha^P}[C] = 1\) perché \(\Val{\alpha^P}[\neg P] = 1\).
  Se invece \(\neg P \not \in C \implies C \setminus \set{P} \in S_P\). Dunque \(\Val{\alpha}[C\setminus\set{\neg P}] = 1 \implies \Val{\alpha^P}[C\setminus\set{P}] = 1\). Ma \(\Val{\alpha}[C\setminus\set{\neg P}] = 1\) vale solo se \(\alpha^P\) soddisfa almeno un letterale di \(C\) che non è \(P\). Tale letterale è però presente anche in \(C\), dunque anche \(\Val{\alpha^P}[C] = 1\).

  Ora possiamo finalmente applicare l'ipotesi induttiva, avendo che \(S_P\models_R \square\) e \(S_{\neg P}\models_R \square \).
  In particolare quindi esistono due alberi di derivazione del calcolo di risoluzione che hanno uno dei due insiemi come foglie e \(\square\) in radice. Chiamiamoli \(\pi_P\) e \(\pi_{\neg P}\) e valutiamo i diversi casi:
  \begin{enumerate}
    \item Se tutte le foglie di \(\pi_P\) o di \(\pi_{\neg P}\) sono clausole di \(S\) per monotonicità la deduzione di \(\square\) è anche una deduzione da \(S\), e dunque \(S\models_R \square\).
    \item Esiste almeno una foglia \(C_P\) di \(\pi_P\) e almeno una \(C_{\neg P}\) di \(\pi_{\neg P}\) con \(C_P,C_{\neg P} \not\in S\).
        In altre parole c'è almeno una foglia per deduzione che è stata ottenuta dalla sottrazione del letterale.
        Consideriamo ora una derivazione per risoluzione da \(C_P\) e una qualsiasi altra clausola in clash \(D\). Essendo che \(P\notin C_P\) e \(P\not\in D\), \(P\) non ci sarà neanche nel loro risolvente \(R\). A questo punto però anche se aggiungessi \(P\) sia a \(C_P\) che a \(R\) avrei comunque un applicazione corretta di risoluzione. Con questa idea possiamo aggiungere \(P\) a tutte le \(C \in S_P\setminus S\). Facendo ciò essendo che \(P\) non poteva essere mai usata per un clash nella deduzione tale \(P\) aggiunta verrà propagata fino alla fine della deduzione, ottenendo quindi una deduzione di \(P\) da \(S\) a partire di quella di \(\square\) da \(S_P\). Applicando il ragionamento duale per \(S_{\neg P}\) otteniamo una deduzione di \(\neg P\) da \(S\). A questo punto è sufficiente un ulitma applicazione di risoluzione per ottenere \(\square\) da \(S\) facendo fare clash alle clausole derivate \(\set{P}\) e \(\set{\neg P}\)
  \end{enumerate}
\end{proof}

\begin{example}{}
  Supponiamo di avere
  \[S = \set{\set{P,Q}, \set{P,\neg Q,\neg R}, \set{R}, \set{\neg P,\neg R}, \set{\neg P,Q,R} }\]
  Fissando \(P\) possiamo creare gli insiemi
  \[S_P=\set{\set{Q},\set{\neg Q, \neg R}, \set{R}}\]
  \[S_{\neg P} = \set{\set{R}, \set{\neg R}, \set{Q,R}}\]

  Ora è vacile vedere che \(S_P \models_R\square\) e \(S_{\neg P} \models_R \square\):
  AGGIUNGI DEDUZIONE

  Per queste deduzioni abbiamo che \(\set{\neg R}\) in \(S_{\neg P}\) non apparteneva a \(S\), così come \(\set{\neg Q, \neg R}\).
  Aggiungiamo quindi alla deduzione \(P\), e propaghiamola.

  AGGIUNGI NUOVA DEDUZIONE

  A questo punto possiamo dedurre \(\square\) facilmente col clash di \(P\) e \(\neg P\).
\end{example}

Come già anticipato questa forma di completezza è più debole di quella vista per la deduzione naturale perché non è completa per la conseguenza logica. Infatti seppur \(A\models_{ND}A\lor B\), \(\set{A}\models_R\set{A,B}\).

L'unica cosa che si può fare è mostrare che \(\set{A,\neg(A\lor B) }\) è indossisfacibile. Questo è facile poiché \(\neg(A\lor B)\equiv \neg A \land \neg B\), quindi l'insieme diventa \(\set{\set{A},\set{\neg A}, \set{B}}\) da cui è facilmente deducibile square.

\section{Applicazioni della logica propisizionale}

Avendo trovato un algoritmo corretto, completo e terminante per risolvere la soddisfacibilità di formule proposizionali, se è possibile riuscire a trovare una riduzione di un qualsiasi problema a \(\Sat\) possiamo avere un algoritmo per risolvere quel problema.

Vediamo quindi qualche problema e qualche riduzione.

Supponiamo di avere un grafo \(G=(V,E)\) orientato e vogiamo vedere se \(G\) è ciclico.
Chiaramente è noto un algoritmo per verificare che un grafo sia ciclico, ma vediamo come potremmo affrontarlo tramite logia proposizionale.

Per fare una riduzione dobbiamo avere un modo di rappresentare l'input e la soluzione. Inoltre, a differenza di un tipo di soluzione operativa tipica di altri modelli di calcolo, in cui delle operazioni fanno del lavoro, la logica ha un approccio prettamente dichiarativo, in cui esprimiamo solo i vincoli che devono sussistere tra l'input e l'output.

In particolare per rappresentare il grafo, andremo a introdurre per ogni coppia \(u,v \in V\) una proposizione atomica \(X_{u,v}\).

A questo punto per rappresentare gli archi è possibile costruire la formula \(\phi_E = \bigwedge_{(u,v)\in E} X_{u,v} \land \bigwedge_{(u,v)\not\in E}\neg X_{u,v}\). Intuitivamente tale formula rappresenta un'assegnamento soddisfatto solo dalla matrice di adiacenza del grafo.

Per l'output vorremo mostrare che esista un \(S\subseteq E\) tale che \(S\) contiene un ciclo. Partiamo da definire le variabili per l'output come \(Y_{u,v}\) per ogni \(u,v \in V\). Dunque possiamo definire \(\phi_S = \bigwedge_{u,v \in V} Y_{u,v}\to X_{u,v}\). Tale formula ci dice che ogni elemento di \(S\) deve stare in \(E\). Inoltre, per far si che \(S\) contenga un ciclo non può essere vuoto, e questo lo posso esprimere con \(\phi_{\neq \emptyset} = \bigvee_{u,v\in V} Y_{u,v} \).

Ci manca solo descrivere che dentro \(S\) ci sia un ciclo.
Per fare ciò useremo la formula \(\phi_{cyc} = \bigwedge_{u,v} (Y_{u,v}\to \bigvee_{z \in V} Y_{v,z})\).
Per far vedere che questa condizione rappresenta la presenza di un ciclo assumiamo che tutte le \(\phi\) viste fino ad ora siano soffisfatte. Essendo che \(\phi_{\neq \emptyset}\), possiamo prendere un qualsiasi \(Y_{u,v}\), essendo che tale proposizione è vera da \(\phi_{cyc}\) dev'esserci un altro arco che parte da \(v\) in \(S\). Iterando il procedimento possiamo vedere che questo porta a 
costruire un percorso infinito, che su un grafo finito implica l'esistenza di un ciclo.

Nell'altro verso, se è presente un ciclo esiste un percorso infinito nel grafo. Allora esiste un assegnamento delle variabili \(Y_{u,v}\), in particolare quello che assegna a \(1\) le variabili corrispondenti agli archi del percorso. Dunque se esiste un percorso infinito tale formula è soddisfatta.

A questo punto il grafo ha un ciclo se e solo se \(\phi_E \land \phi_S \land \phi_{\neq \emptyset} \land \phi_{cyc} \) è soddisfacibile.

Vediamo ora un problema più complesso.

Sia \(G = (V,E)\) un grafo e un insieme \(C = \set{c_1, \ldots, c_k}\) di colori.
Il problema di \(k-\)colorabilità di un grafo si chiede se \(\exists f: V\to C \st \forall u,v \in V, (u,v) \in E \implies f(u)\neq f(v)\).

Questo problema è molto interessante perché rappresenta una generalizzazione di un problema di allocazione di risorse.

Questo problema è molto più complicato del precedente perché si dimostra essere \(NP-\)completo.

Vediamo come ridurre tale problema in logica. La funzione di output la possiamo vedere come una relazione \(f\subseteq V\times C\) con un vincolo funzionale. Dunque possiamo rappresentarla estensionalmente tramite le coppie del prodotto cartesiano.

In particolare, possiamo avere \(X_v^c\) con \(v\in V\) e \(c\in C\), in cui ogni assegnamento di tali variabili va interpretato come una colorazione del vertice \(v\) del colore \(c\) se \(X_v^c\) vale \(1\). Dunque \(\AP = \set{X_v^c}[c\in C \land v \in V]\) e quindi \(\abs{\AP} = \abs{V}k\).

Dobbiamo però imporre il vincolo di funzionalità e di totalità e quello dell'adiacenza con colori diversi.

Il primo, \(\phi_f = \bigwedge_{v\in V}(\bigvee_{c\in C} X_v^c \land \bigvee_{c'\in C\setminus\set{c}} \neg X_v^{c'})\). Questo mi dice che ogni vertice deve avere un colore, e che non ha nessuno degli altri. La dimensione di questa formula è \(\abs{V}k^2\)

Il secondo è molto semplice, poiché basta definire \(\phi_2 = \bigwedge_{c \in C}\bigwedge_{u,v\in V \st (u,v)\in E} (\neg (X_u^c \land X_v^c))\). La lunghezza di questa formula è \(\abs{V}^2k\).

Dunque \(\phi_f \land \phi_2 = \BigO(\abs{V}^3)\). Questa codifica del problema è però problematica se l'input dei colori lo rappresento solo usando \(k\). Questo perché un numero per essere rappresentato necessita \(\log_2(k)\) celle. Quindi avere una formula di dimensione polinomiale in \(k\) non è una formula di dimensione polinomiale in \(\lceil\log_2(k)\rceil\).

Esiste però una codifica più efficiente. Visto che \(k\) può essere espresso con \(\lceil\log_2(k)\rceil\), posso rappresentare i colori con \(\lceil\log_2(k)\rceil\) variabili che rappresentano i bit. A questo punto per rappresentare ogni vertice utilizzo le variabili \(X_0^v\ldots X_{\log_2(k)-1}^v\) il cui assegnamento rappresenta il colore assegnato. Questo fa in modo di avere anche il vincolo funzionale, perché ogni assegnamento può assegnare a vero solo un colore a ogni vertici. Abbiamo chiaramente ora che \(\abs{AP} = \abs{V}\log_2(k)\).

Dobbiamo solo esprimere il vincolo di adiacenza con colori diversi. Possiamo farlo con
\[\phi_2' = \land_{u,v \in V \st (u,v)\in E}\left( \bigvee_{i=0}^{\log_2(k)-1} (X_i^u \land \neg X_i^v) \lor (\neg X_i^u \land X_i^v)\right)\] 

Tale codifica chiaramente occupa spazio \(\abs{V}^2\log_2(k)\), che è lineare nell'input.

Possiamo notare che in questa codifica, per \(k=2\) otteniamo un problema in \(2-\Sat\), il che mostra facilmente come il 2 graph coloring sia polinomiale.